\documentclass[10pt,a4paper]{amsart}
\usepackage[margin=19mm]{geometry}
\usepackage{amsmath,amssymb,mathtools}
\usepackage[scr=rsfs,cal=euler]{mathalfa}
\usepackage{xfrac}
\usepackage[unicode,hidelinks,bookmarks=false]{hyperref}
\newcommand{\ie}{i.e.\ }
\newcommand{\cf}{cf.\ }
\newcommand{\resp}{respectively\ }
\newcommand{\Subsection}{\subsection}
\providecommand{\nobreakdash}{\nobreak\hbox{-}}
\setlength{\emergencystretch}{2em}
\multlinegap0pt
\usepackage{amssymb}
\usepackage{natbib}
\newtheorem{theorem}{Theorem}
\newtheorem{lemma}{Lemma}
\title{\textbf{Fractional Grönwall--Wendroff Inequalities for Implicit Systems with Distributed Memory}}
\author{R\^omulo Damasclin Chaves dos Santos}
\date{Source: arXiv:2602.06978v1 (CC BY 4.0)}
\begin{document}
\maketitle

\noindent\emph{Source and licence.} \textbf{Fractional Grönwall--Wendroff Inequalities for Implicit Systems with Distributed Memory}. Version arXiv:2602.06978v1 (CC BY 4.0), distributed by arXiv under the Creative Commons Attribution 4.0 International licence, http://creativecommons.org/licenses/by/4.0/, as recorded in the arXiv licence metadata for this exact version and shown on the abstract page https://arxiv.org/abs/2602.06978v1. Full licence notice: \texttt{licenses/arxiv-2602.06978-CC-BY-4.0.txt}.\par\smallskip

\noindent\textit{Editorial scope.} Counted unit: the article's theorem thm:limit\_cycle (source0.tex lines 368--393). The article's complete original proof of it is delivered with it (source0.tex lines 395--482, one \texttt{proof} environment of 88 lines). No mathematics is written, repaired or re-derived: the statement and the proof are the article's own lines, in the article's own notation, and no retained line is substituted or re-flowed.  Retained uncounted as the source-local context the proof needs: lines 110--122; lines 287--299; lines 315--322.  Nothing was omitted from any retained span, so the package carries no omission record.  No retained line contains a citation, so the wrapper carries no bibliography and no bibliography entry.  No retained line contains a figure environment, an image-inclusion command or drawing code, so no image is delivered and the package declares no asset dependency.  Editorial formatting only, in addition: an AMS \texttt{amsart} preamble that restates the article's own declarations and macros used by the retained lines, namely the environments \texttt{theorem}, \texttt{lemma} on separate counters, together with the standard packages those lines need.  The retained environments print in this document as Theorem 1, Lemma 1 in order of appearance, whereas the article prints the same items with the numbering of its own full document.  The complete native source of the article as distributed in the arXiv e-print for this exact version is bundled unchanged as \texttt{sources/194144/source0.tex}.
	\begin{theorem}[Multivariate Fractional Grönwall--Wendroff Inequality]
		\label{thm:main_ineq}
		Let $u,f \in C([0,T],\mathbb{R}^n)$ with $u(t)\ge 0$ and $f(t)\ge 0$ componentwise. Assume that
		\begin{equation}
			u(t) \le f(t) + \int_0^t (t-s)^{\alpha-1} A(s)u(s)\,ds + \int_0^t (t-s)^{\beta-1} B(s)u(s-\tau)\,ds,
			\label{eq:inequality_form}
		\end{equation}
		for all $t\in[0,T]$, where $\alpha,\beta\in(0,1]$, $A,B\in L^\infty([0,T],\mathbb{R}^{n\times n})$, and $u(t)=\phi(t)$ for $t\in[-\tau,0]$. Then there exists a constant $M>0$, depending only on $\alpha,\beta,\|A\|_\infty,\|B\|_\infty$ and $T$, such that
		\begin{equation}
			u(t) \le M\Bigl(\|\phi\|_{[-\tau,0]}+\sup_{s\in[0,t]}f(s)\Bigr), \qquad \forall\, t\in[0,T].
			\label{eq:inequality_result}
		\end{equation}
	\end{theorem}

	The system is given by
	\begin{equation}
		\begin{cases}
			^{C}\mathcal{D}_{0}^{\alpha} v(t) = v(t) - \dfrac{v^3(t)}{3} - w(t) + I_{\mathrm{ext}} + \lambda v(t-\tau), \\[0.3em]
			^{C}\mathcal{D}_{0}^{\alpha} w(t) = \varepsilon \bigl(v(t) + a - b w(t)\bigr),
		\end{cases}
		\label{eq:fhn_model}
	\end{equation}
	where $0<\alpha<1$, $\varepsilon>0$ is a small parameter, $\lambda\in\mathbb{R}$ denotes the feedback intensity, and $\tau>0$ is the delay. The system is endowed with functional initial conditions
	\begin{equation}
		v(t)=\phi_v(t), \qquad w(t)=\phi_w(t), \qquad t\in[-\tau,0].
		\label{eq:fhn_initial}
	\end{equation}

	\begin{lemma}[Fractional dissipativity]
		\label{lem:dissipativity}
		There exist constants $R>0$ and $T_0>0$ such that every solution of \eqref{eq:fhn_model} satisfies
		\begin{equation}
			\|\mathbf{x}(t)\| = \sqrt{v^2(t)+w^2(t)} \le R, \qquad \forall\, t\ge T_0.
			\label{eq:absorbing_ball}
		\end{equation}
	\end{lemma}

\begin{theorem}[Existence and Stability of Fractional Limit Cycles]
	\label{thm:limit_cycle}
	Consider the delayed fractional FitzHugh--Nagumo system \eqref{eq:fhn_model} with $0<\alpha<1$. 
	Assume the following conditions hold:
	\begin{enumerate}
		\item \textbf{Small delay coupling:} $|\lambda| < \lambda_0(\alpha, \tau)$, where
		\begin{equation}
			\lambda_0(\alpha, \tau) = \frac{\Gamma(\alpha+1)}{\tau^\alpha} \cdot \min\left\{\frac{1}{2}, \frac{1-\alpha}{\alpha}\right\};
			\label{eq:lambda_condition}
		\end{equation}
		
		\item \textbf{Slow recovery dynamics:} $0 < \varepsilon < \varepsilon_0(\alpha, b)$, where
		\begin{equation}
			\varepsilon_0(\alpha, b) = \frac{b \Gamma(\alpha+1)}{2} \left(1 - \frac{|\lambda|\tau^\alpha}{\Gamma(\alpha+1)}\right);
			\label{eq:epsilon_condition}
		\end{equation}
		
		\item \textbf{Subthreshold parameters:} $a$ and $b$ satisfy $a < 1 + \frac{b^2}{4\varepsilon}$.
	\end{enumerate}
	Then there exists at least one nontrivial periodic solution $\mathbf{x}^*(t)$ of \eqref{eq:fhn_model} with period $T^* > 0$. Moreover, this solution lies in an invariant annular region $\mathcal{A} \subset \mathbb{R}^2$ defined by
	\begin{equation}
		\mathcal{A} = \left\{ (v,w) \in \mathbb{R}^2 : R_1 \leq \sqrt{v^2 + \frac{w^2}{\varepsilon}} \leq R_2 \right\},
		\label{eq:annular_region}
	\end{equation}
	where $0 < R_1 < R_2 < \infty$ depend explicitly on $\alpha, \tau, \lambda, \varepsilon, a, b$.
\end{theorem}

\begin{proof}
	We present a detailed proof based on Poincaré maps, dissipativity estimates, and fixed-point theory.
	
	\noindent\textbf{1.~Construction of the Poincaré map.}
	Let $\mathcal{X} = C([-\tau,0], \mathbb{R}^2)$ equipped with the supremum norm $\|\phi\|_{\mathcal{X}} = \sup_{\theta \in [-\tau,0]} |\phi(\theta)|$. 
	For $\phi \in \mathcal{X}$, denote by $\mathbf{x}(t;\phi) = (v(t;\phi), w(t;\phi))^\top$ the unique solution of \eqref{eq:fhn_model} with initial history $\phi$.
	
	Define the time-$T$ Poincaré map $P_T: \mathcal{X} \to \mathcal{X}$ by
	\begin{equation}
		(P_T\phi)(\theta) = \mathbf{x}(\theta + T; \phi), \quad \theta \in [-\tau, 0],
		\label{eq:poincare_map_T}
	\end{equation}
	where $T > \tau$ will be chosen appropriately.
	
	\noindent\textbf{2.~Invariant annular region.}
	From Lemma~\ref{lem:dissipativity}, we have the existence of absorbing balls. 
	A refined analysis using the Lyapunov functional $V(t) = \frac{1}{2}v^2(t) + \frac{1}{2\varepsilon}w^2(t)$ yields
	\begin{equation}
		^{C}\mathcal{D}_0^\alpha V(t) \leq -\delta V(t) + C_1 V(t-\tau) + C_2,
		\label{eq:lyapunov_inequality}
	\end{equation}
	with $\delta = \min\left\{\frac{2}{3}, \frac{b}{2}\right\}$, $C_1 = \frac{|\lambda|}{2}$, and $C_2 = \frac{|I_{\text{ext}}|^2}{2} + \frac{\varepsilon a^2}{2}$.
	
	Applying the fractional Grönwall inequality (Theorem~\ref{thm:main_ineq}) to \eqref{eq:lyapunov_inequality}, we obtain
	\begin{equation}
		V(t) \leq M\left(V(0) + \frac{C_2}{\delta}\right) E_\alpha\left(-\frac{\delta}{2} t^\alpha\right) + \frac{2C_2}{\delta},
		\label{eq:lyapunov_bound}
	\end{equation}
	where $E_\alpha(z) = \sum_{k=0}^\infty \frac{z^k}{\Gamma(\alpha k + 1)}$ is the Mittag-Leffler function and $M > 0$ depends on $\alpha, \tau, \lambda$.
	
	This implies that all solutions eventually enter and remain in the annular region $\mathcal{A}$ defined in \eqref{eq:annular_region} with
	\begin{equation}
		R_1 = \sqrt{\frac{C_2}{\delta}}, \quad R_2 = \sqrt{\frac{2C_2}{\delta} + \frac{M C_2}{\delta}}.
		\label{eq:radius_bounds}
	\end{equation}
	
	\noindent\textbf{3.~Compactness of the Poincaré map.}
	Let $\mathcal{B} = \{\phi \in \mathcal{X} : \|\phi\|_{\mathcal{X}} \leq R_2\}$.
	From the integral representation of solutions
	\begin{equation}
		\mathbf{x}(t) = \mathbf{x}(0) + \frac{1}{\Gamma(\alpha)} \int_0^t (t-s)^{\alpha-1} \mathbf{F}(s,\mathbf{x}(s),\mathbf{x}(s-\tau))\, ds,
		\label{eq:solution_representation}
	\end{equation}
	we derive Hölder regularity estimates. For $t_1, t_2 \in [0, T]$ with $|t_1 - t_2| < \delta$,
	\begin{equation}
		|\mathbf{x}(t_1) - \mathbf{x}(t_2)| \leq \frac{L}{\Gamma(\alpha+1)} |t_1 - t_2|^\alpha,
		\label{eq:holder_estimate}
	\end{equation}
	where $L$ depends on $R_2$ and the Lipschitz constant of $\mathbf{F}$.
	
	By the Arzelà--Ascoli theorem, $P_T(\mathcal{B})$ is relatively compact in $\mathcal{X}$.
	Moreover, from \eqref{eq:lyapunov_bound}, $P_T(\mathcal{B}) \subseteq \mathcal{B}$ for sufficiently large $T$.
	
	\noindent\textbf{4.~Existence of a fixed point.}
	Choose $T > \max\left\{\tau, \left(\frac{2}{\delta}\ln\left(\frac{2M}{\varepsilon}\right)\right)^{1/\alpha}\right\}$ to ensure
	\begin{equation}
		E_\alpha\left(-\frac{\delta}{2} T^\alpha\right) < \frac{\varepsilon}{2M}.
		\label{eq:period_condition}
	\end{equation}
	
	The Poincaré map $P_T: \mathcal{B} \to \mathcal{B}$ is continuous (by continuous dependence on initial conditions) and compact. 
	Applying Schauder's fixed-point theorem, there exists $\phi^* \in \mathcal{B}$ such that
	\begin{equation}
		P_T \phi^* = \phi^*.
		\label{eq:fixed_point}
	\end{equation}
	
	\noindent\textbf{5.~Nontriviality and periodicity.}
	The corresponding solution $\mathbf{x}^*(t) = \mathbf{x}(t; \phi^*)$ satisfies
	\begin{equation}
		\mathbf{x}^*(t+T) = \mathbf{x}^*(t) \quad \forall t \geq 0,
		\label{eq:periodicity}
	\end{equation}
	with period $T^*$ dividing $T$. To show nontriviality, we verify that $\mathbf{x}^*(t)$ is not an equilibrium. 
	For $|\lambda|$ sufficiently small and $\varepsilon$ satisfying \eqref{eq:epsilon_condition}, linear stability analysis shows that the unique equilibrium $(v_0, w_0)$ with $w_0 = (v_0+a)/b$ and $v_0$ solving
	\begin{equation}
		v_0 - \frac{v_0^3}{3} - \frac{v_0+a}{b} + I_{\text{ext}} + \lambda v_0 = 0
		\label{eq:equilibrium}
	\end{equation}
	is unstable for $\alpha < 1$. The instability follows from the characteristic equation
	\begin{equation}
		s^\alpha = 1 - v_0^2 - \frac{1}{b} + \lambda e^{-s\tau},
		\label{eq:characteristic}
	\end{equation}
	which admits solutions with $\Re(s) > 0$ under conditions \eqref{eq:lambda_condition} and \eqref{eq:epsilon_condition}.
	
	Therefore, $\mathbf{x}^*(t)$ is a nontrivial periodic solution, i.e., a limit cycle of the FHN-$\alpha$-$\tau$ system.
\end{proof}

\end{document}
